\section{Generating Random Numbers}

% \subsection{Random Number Generation}

% \subsubsection*{Uniform inputs}

% Monte Carlo simulation is driven by apparently random numbers. Ideally we want $U_1,U_2,\ldots$ such that each $U_i\sim\mathrm{Unif}[0,1]$ and the $U_i$ are mutually independent. Independence is crucial: $U_i$ should not be predictable from $U_1,\ldots,U_{i-1}$.

% A pseudorandom number generator produces a finite deterministic sequence $u_1,\ldots,u_K\in[0,1]$. A good generator makes short segments hard to distinguish from independent uniforms. Desirable properties: long period, reproducibility, speed, portability, and statistical randomness.

% \subsubsection*{Linear congruential generator}

% A linear congruential generator is

% $x_{i+1}=a x_i \bmod m,\qquad u_{i+1}=x_{i+1}/m$.

% Here $a$ is the multiplier, $m$ the modulus, and $x_0\in\{1,\ldots,m-1\}$ the seed. The modulo operation is

% $y\bmod m=y-\lfloor y/m\rfloor m$.

% Once a value repeats, the whole sequence repeats. The generator has full period if it produces all $m-1$ nonzero values before repeating. Large $m$ alone is not enough; poor choices of $a,m$ can produce short cycles.

\subsection{General Sampling Methods}

\subsubsection*{Uniform distribution}

Assume independent uniforms satisfying $P(U_i\le u)=0$ if $u<0$, $P(U_i\le u)=u$ if $0\le u\le1$, and $P(U_i\le u)=1$ if $u>1$. Sampling methods transform these uniforms into variables with desired distributions.

\subsection{Inverse Transform Method}

\subsubsection*{General method}

To sample $X$ with CDF $F$, set

$X=F^{-1}(U),\qquad U\sim\mathrm{Unif}[0,1]$.

If $F$ is not strictly increasing, use the generalized inverse

$F^{-1}(u)=\inf\{x:F(x)\ge u\}$.

This chooses the smallest compatible $x$. Validity follows from

$P(X\le x)=P(F^{-1}(U)\le x)=P(U\le F(x))=F(x)$.

Flat sections of $F$ have zero probability: if $F$ is constant on $[a,b]$, then $P(a<X\le b)=F(b)-F(a)=0$.

\subsubsection*{Examples}

Exponential with mean $\theta$: $F(x)=1-e^{-x/\theta}$ for $x\ge0$, so

$X=-\theta\log(1-U)$, equivalently $X=-\theta\log U$.

Rayleigh example:
$F(x)=1-e^{-2x(x-b)}, x\ge b$,

so the inverse-transform sample is
$X=\frac{b}{2}+\frac{\sqrt{b^2-2\log U}}{2}$.

Discrete distribution: for outcomes $c_1<\cdots<c_n$ with probabilities $p_i$, define $q_0=0$ and
$q_i=\sum_{j=1}^i p_j$.
Generate $U$, find $K$ with $q_{K-1}<U\le q_K$, and set $X=c_K$. Binary search can implement the lookup.

\subsubsection*{Conditional sampling}

To sample $X$ conditional on $a<X\le b$, with $F(a)<F(b)$, generate

$V=F(a)+[F(b)-F(a)]U$.

Then $F^{-1}(V)$ has the conditional distribution because

$P(F^{-1}(V)\le x)=\frac{F(x)-F(a)}{F(b)-F(a)}$.

\subsection{Acceptance-Rejection Method}

\subsubsection*{General method}

To sample from a target density $f$ on $\mathcal{X}$, choose a convenient proposal density $g$ and a constant $c$ such that
$f(x)\le c g(x),\qquad x\in\mathcal{X}$.

Generate $X\sim g$ and $U\sim\mathrm{Unif}[0,1]$. Accept $X$ if
$U\le\frac{f(X)}{c g(X)}$;

otherwise reject and repeat. If $Y$ is the accepted value, then

$P(Y\in A)=P(X\in A\mid U\le f(X)/(cg(X)))$.

The acceptance probability is

$P(U\le f(X)/(cg(X)))=\int_{\mathcal{X}}\frac{f(x)}{cg(x)}g(x)\,dx=1/c$.

Hence
$P(Y\in A)=c\int_A\frac{f(x)}{cg(x)}g(x)\,dx=\int_A f(x)\,dx$.

Thus $Y$ has density $f$. The number of proposals until acceptance is geometric with mean $c$, so smaller $c$ and tighter bounds are better.

% \subsubsection*{Beta distribution}

% The beta density on $[0,1]$ is

% $f(x)=\frac{1}{B(\alpha_1,\alpha_2)}x^{\alpha_1-1}(1-x)^{\alpha_2-1}$.

% Here

% $B(\alpha_1,\alpha_2)=\int_0^1x^{\alpha_1-1}(1-x)^{\alpha_2-1}\,dx=\frac{\Gamma(\alpha_1)\Gamma(\alpha_2)}{\Gamma(\alpha_1+\alpha_2)}$,

% and

% $\Gamma(z)=\int_0^\infty x^{z-1}e^{-x}\,dx,\qquad \operatorname{Re}(z)>0$.

% For integer $n>0$, $\Gamma(n)=(n-1)!$. If $\alpha_1,\alpha_2\ge1$ and at least one exceeds $1$, the mode is

% $\frac{\alpha_1-1}{\alpha_1+\alpha_2-2}$.

% Let $c$ be the maximum value of $f$ and use $g(x)=1$ on $[0,1]$. Generate $U_1,U_2$ until $cU_2\le f(U_1)$, then return $U_1$.

\subsubsection*{Normal from double exponential}

For the double exponential density $g(x)=e^{-|x|}/2$ and standard normal density

$f(x)=\frac{1}{\sqrt{2\pi}}\exp(-x^2/2)$,

the ratio satisfies
$\frac{f(x)}{g(x)}=\sqrt{\frac{2}{\pi}}e^{-x^2/2+|x|}\le\sqrt{\frac{2e}{\pi}}\approx1.3155\equiv c$.

The rejection test may be written as

$u>\exp\left(-\frac12x^2+|x|-\frac12\right)=\exp\left(-\frac12(|x|-1)^2\right)$.

Using symmetry: generate $U_1,U_2,U_3$, set $X=-\log U_1$, reject if $U_2>\exp[-0.5(X-1)^2]$, then set $X\leftarrow -X$ if $U_3\le0.5$.

\subsection{Normal Random Variables and Vectors}

\subsubsection*{Multivariate normal}

A $d$-dimensional normal is written $N(\mu,\Sigma)$, where $\mu$ is a $d$-vector and $\Sigma$ is a $d\times d$ covariance matrix. A covariance matrix must be symmetric and positive semidefinite:
$x^\top\Sigma x\ge0,\qquad x\in\mathbb{R}^d$.

If $\Sigma$ is positive definite, the density is

$\phi_{\mu,\Sigma}(x)=\frac{1}{(2\pi)^{d/2}|\Sigma|^{1/2}}\exp\left[-\frac12(x-\mu)^\top\Sigma^{-1}(x-\mu)\right]$.

If $X\sim N(\mu,\Sigma)$, then $X_i\sim N(\mu_i,\sigma_i^2)$ with $\sigma_i^2=\Sigma_{ii}$, and

$\operatorname{Cov}(X_i,X_j)=\mathbb{E}[(X_i-\mu_i)(X_j-\mu_j)]=\Sigma_{ij}$,
$\rho_{ij}=\frac{\Sigma_{ij}}{\sigma_i\sigma_j}$, $\Sigma_{ij}=\sigma_i\sigma_j\rho_{ij}$.

If $\Sigma$ is positive semidefinite but singular, define $X=\mu+AZ$ with $Z\sim N(0,I_d)$ and $AA^\top=\Sigma$. Then $X$ has no full $d$-dimensional density when $\operatorname{rank}(\Sigma)<d$.

\subsubsection*{Useful properties}

Linear transformation: if $X\sim N(\mu,\Sigma)$, then
$AX\sim N(A\mu,A\Sigma A^\top)$.

Conditioning formula: for a partitioned normal vector with $\Sigma_{[22]}$ full rank,

$(X_{[1]}\mid X_{[2]}=x)\sim N(\mu_{[1]}+\Sigma_{[12]}\Sigma_{[22]}^{-1}(x-\mu_{[2]}),\ \Sigma_{[11]}-\Sigma_{[12]}\Sigma_{[22]}^{-1}\Sigma_{[21]})$.

Moment generating function:
$\mathbb{E}[\exp(\theta^\top X)]=\exp(\mu^\top\theta+\frac12\theta^\top\Sigma\theta)$.

\subsection{Generating Normals}

\subsubsection*{Box-Muller method}

The Box-Muller method generates a bivariate standard normal. If $Z\sim N(0,I_2)$, then $R=Z_1^2+Z_2^2$ is exponential with mean $2$:
$P(R\le x)=1-e^{-x/2}$.

Given $R$, $(Z_1,Z_2)$ is uniformly distributed on the circle of radius $\sqrt{R}$. Generate independent $U_1,U_2$, set
$R=-2\log U_1,\qquad V=2\pi U_2$,

and return
$Z_1=\sqrt{R}\cos V,\qquad Z_2=\sqrt{R}\sin V$.

\subsubsection*{Multivariate normals}

To sample $X\sim N(\mu,\Sigma)$, generate $Z\sim N(0,I_d)$, find $A$ such that
$AA^\top=\Sigma$,

and set $X=\mu+AZ$. Then $X\sim N(\mu,\Sigma)$.

\subsection{Cholesky Factorization}

\subsubsection*{Positive definite case}

Cholesky represents $\Sigma=AA^\top$ with $A$ lower triangular. For

$\Sigma=\begin{pmatrix}\sigma_1^2&\sigma_1\sigma_2\rho\\ \sigma_1\sigma_2\rho&\sigma_2^2\end{pmatrix}$,

a Cholesky factor is
$A=\begin{pmatrix}\sigma_1&0\\ \rho\sigma_2&\sqrt{1-\rho^2}\sigma_2\end{pmatrix}$.

Thus generate
$X_1=\mu_1+\sigma_1Z_1$,
$X_2=\mu_2+\sigma_2\rho Z_1+\sigma_2\sqrt{1-\rho^2}Z_2$.

For general $d$, use
$\Sigma_{ij}=\sum_{k=1}^jA_{ik}A_{jk},\qquad j\le i$.

The recursive formulas are

$A_{ij}=\frac{\Sigma_{ij}-\sum_{k=1}^{j-1}A_{ik}A_{jk}}{A_{jj}},\, j<i$,
$\quad A_{ii}=\sqrt{\Sigma_{ii}-\sum_{k=1}^{i-1}A_{ik}^2}$.
Loop over columns $j=1,\ldots,d$ and rows $i=j,\ldots,d$, solving one new entry at a time.

\subsubsection*{Semidefinite case}

If $\Sigma$ is only positive semidefinite, $\operatorname{rank}(\Sigma)<d$ and some diagonal entry of a lower triangular factor must be zero, so Cholesky may divide by zero or become round-off sensitive. A mathematical fix sets the whole $j$th column of $A$ to zero when $A_{jj}=0$, but exact zero checks are fragile.

Prefer reducing to full rank: choose a subvector $\tilde{X}=(X_{i_1},\ldots,X_{i_k})$ with full-rank covariance $\tilde{\Sigma}$, find $\tilde{A}\tilde{A}^\top=\tilde{\Sigma}$, and sample
$X=D\tilde{A}Z,\qquad Z\sim N(0,I_k)$.
